%% Appendix body, shared by two parents: appendix.tex (the standalone
%% supplementary document shipped in the artifact — POPL'27 submission
%% flow) and ../arxiv.tex (which inlines it after the bibliography).
%% Cross-references to the paper are FROZEN TEXT (Theorem~3.7 etc.), not
%% \Cref: the standalone document cannot resolve the paper's labels, and
%% the numbering is identical in the arXiv version, where the body is
%% inlined unchanged. Crosswalk note in appendix.tex.

\section{Proofs}
\label{app:proofs}

\subsection{Pasting (Theorem~3.7)}

We spell out the chain, including the domain side conditions elided in
the paper's proof sketch. Let $p \in \dom(P_2 \circ P_1)$, i.e.\
$p \in \dom(I_A) \cap \dom(T_1)$, $q := T_1(p) \in \dom(T_2)$,
$r := T_2(q)$, $I_C(r)$ defined,
$\carry_2(I_C(r)) \in \dom(\carry_1)$, and (from $p \in \dom(P_1)$,
$q \in \dom(P_2)$) $I_B(q)$ defined with
$I_B(q) \in \dom(\carry_1)$ and $I_C(r) \in \dom(\carry_2)$.

\begin{enumerate}
\item By faithfulness of $P_1$ at $p$ and $\proj \subseteq \proj_1$,
  $\proj(I_A(p)) = \proj(\carry_1(I_B(q)))$.
\item By faithfulness of $P_2$ at $q$:
  $\proj_2(I_B(q)) = \proj_2(\carry_2(I_C(r)))$.
\item Both $I_B(q)$ and $\carry_2(I_C(r))$ lie in $\dom(\carry_1)$ (the
  first by $p \in \dom(P_1)$, the second by
  $p \in \dom(P_2 \circ P_1)$). By
  $(\proj_2 \Rightarrow \proj)$-support applied to step (2):
  $\proj(\carry_1(I_B(q))) = \proj(\carry_1(\carry_2(I_C(r))))$.
\item Chaining (1) and (3):
  $\proj(I_A(p)) = \proj(\carry_1(\carry_2(I_C(r))))$, which is
  faithfulness of $P_2 \circ P_1$ at $p$ w.r.t.\ $\proj$. \qed
\end{enumerate}

\emph{Necessity of support.} Without the support condition
(Definition~3.6) the theorem is
false: let $F_B = \{g, h\}$, $\proj_2 = \{g\}$, and let $\carry_1$ copy
field $h$ into a $\proj_1$-kept source field. Take $P_2$ faithful (it
preserves $g$) but with $\carry_2(I_C(r))$ differing from $I_B(q)$ on
$h$ --- permitted, since $h \notin \proj_2$. The outer rectangle then
disagrees on the copied field although both inner squares commute.

\subsection{Composed keep-set, and step granularity}

For fieldwise, step-aligned $\carry_1$ (each kept source field $f$
computed per step as $f = \phi_f(\{\,b.g \mid g \in \mathit{deps}(f)\,\})$),
$(\proj_2 \Rightarrow \proj)$-support holds iff
$\mathit{deps}(f) \subseteq \proj_2$ for every $f \in \proj$; hence the
largest admissible $\proj$ is
$\keep = \{ f \in \proj_1 \mid \mathit{deps}(f) \subseteq \proj_2 \}$ and
support for it is checkable syntactically from the dependency sets.

When $\carry_1$ retimes --- one source step corresponds to a window of
target steps, as in C statements over ISA instructions --- model
$\carry_1$ as a monotone window map $w$ from source step indices to
intervals of target step indices plus a fieldwise read within each
window. Support then requires $\mathit{deps}(f) \subseteq \proj_2$ for
reads \emph{anywhere in the window}, and additionally that the window map
itself be determined by $\proj_2$-observables (e.g.\ by a kept program
counter), since window boundaries influence the projected output. Both
conditions are again syntactic given $w$ and the dependency sets.

\subsection{Localization (Corollary~3.8)}

Contrapositive of the pasting theorem (Theorem~3.7) for two hops. For
$n$ hops, induct
on the route: if the outer rectangle over hops $1..n$ fails at $p$ but
hop $1$'s square passes at $p$, then by the two-hop case applied to
(hop 1, rectangle over hops $2..n$) the rectangle over $2..n$ fails at
$T_1(p)$; recurse. Termination yields a least failing hop $i$; within
it, the square oracle's comparison yields the least failing step and a
witnessing field. \qed

\subsection{Weakest link (Proposition~4.2)}

Soundness: if every hop is $\mathsf{universal}$ on its fragment, iterate
the pasting theorem (Theorem~3.7) over all $p$ in the composite
fragment. If some hop is
$\mathsf{perrun}$, the same iteration is available exactly at programs
whose per-hop images passed the oracle --- the per-run set of the route.
If some hop is only $\mathsf{replay}$, purity composes
(Proposition~3.9) but no faithfulness statement is available for the
route beyond it. Optimality: a route with a $\mathsf{perrun}$ hop entails
no universal statement (take a defect outside the checked set); a route
with a $\mathsf{replay}$ hop entails no faithfulness at all (take a pure
but wrong translator). \qed

\subsection{Ratchet (Proposition~4.7)}

Let $I' \sqsupseteq I$ and let $p \in \dom(I)$. Every quantity in
the faithfulness square (Definition~3.5) evaluated at $p$ reads only
$I$-values on
$\dom(I)$-points, which $I'$ preserves; hence the verdict at $p$ is
unchanged. Coverage monotonicity: $\cov$'s per-construct conjunction can
only gain (new constructs enter $\dom$) and never lose (old verdicts
preserved). Composition of extensions is an extension. For the failure
half: a non-extension update changes some $I(p)$, which can flip the
verdict at $p$ for any pair whose square consumed it --- justifying
the versioned-event rule. \qed

\section{Construct inventories and probe corpora}
\label{app:inventories}

Per-language construct inventories (the yardsticks of the paper's
Definition~4.6) and the probe corpora behind every coverage figure
in the paper's \S6 are enumerated in the artifact --- the per-language
inventory modules (\texttt{gurdy/languages/*/inventory.py}), the
per-pair inventories (\texttt{gurdy/pairs/*/inventory.py}), and the
measured coverage data (\texttt{paper/results/data/}); they are
spec-derived: the RISC-V RV64IMC opcode
inventory, the full 144-opcode EVM inventory, the 256-opcode eBPF
inventory slice, the Wasm numeric/parametric/control inventory, the
OpenSMILES construct list, the Python-subset AST allow-list, and the
two hubs' operator inventories.
